%% kapa: exam=1 round=34 year=2023 period=2 form=A subject=accounting subject_ko=회계학 seq=5 number=45
%% source: Q-Net 감정평가사 기출문제 (cst00309) artlSeq=5211789 file=raw/1008/5211789/2023년도 제34회 감정평가사 1차시험 2교시.pdf
%% conversion: kordoc 4.24.51 -> tools/gen1.py (auto); review=proofread-sonnet; flags=none
\begin{question}{45}
재무보고를 위한 개념체계에 관한 설명으로 옳지 않은 것은?
\begin{choices}
  \item 개념체계는 특정 거래나 다른 사건에 적용할 회계기준이 없는 경우에 재무제표 작성자가 일관된 회계정책을 개발하는 데 도움을 준다.
  \item 개념체계의 어떠한 내용도 회계기준이나 회계기준의 요구사항에 우선하지 아니한다.
  \item 일반목적재무보고의 목적을 달성하기 위해 회계기준위원회는 개념체계의 관점에서 벗어난 요구사항을 정하는 경우가 있을 수 있다.
  \item 개념체계는 수시로 개정될 수 있으며, 개념체계가 개정되면 자동으로 회계기준이 개정된다.
  \item 개념체계에 기반한 회계기준은 경영진의 책임을 묻기 위한 필요한 정보를 제공한다.
\end{choices}
\end{question}
\answer{4}
